 زبان $A_{WOTM}$ {تصمیم‌ناپذیر} است.
برای اثبات این موضوع، نشان می‌دهیم که یک \lr{WOTM} از نظر قدرت محاسباتی دقیقا هم‌ارز یک ماشین تورینگ استاندارد است. با توجه به این هم‌ارزی، مسئله پذیرش برای آن به مسئله پذیرش ماشین تورینگ استاندارد (یعنی $A_{TM}$) کاهش می‌یابد و چون می‌دانیم $A_{TM}$ تصمیم‌ناپذیر است، $A_{WOTM}$ نیز تصمیم‌ناپذیر خواهد بود.

نحوه شبیه‌سازی یک ماشین تورینگ استاندارد مانند $M$ توسط یک \lr{WOTM} مانند $W$ به شرح زیر است:
\begin{itemize}
    \item ماشین $W$ نمی‌تواند نماد خانه‌های پر شده را تغییر دهد، بنابراین به جای بازنویسی روی خانه‌های قبلی، تاریخچه‌ای از تمام پیکربندی‌های ماشین $M$ را روی نوار خود می‌نویسد. هر پیکربندی با یک نماد جداکننده مانند $\#$ از پیکربندی بعدی جدا می‌شود.
    
    \item برای شبیه‌سازی هر مرحله از کار ماشین $M$، ماشین $W$ باید پیکربندی فعلی را در اولین فضای خالی سمت راست نوار کپی کند و تغییرات لازم (نوشتن نماد جدید، تغییر حالت و حرکت کلاهک) را در حین این کپی کردن روی پیکربندی جدید اعمال کند.
    
    \item چالشی که وجود دارد این است که $W$ برای کپی کردن نمادها نمی‌تواند روی آن‌ها علامت‌گذاری کند (زیرا حق تغییر نمادهای غیرخالی را ندارد). اما می‌تواند بدون نیاز به علامت‌گذاری، عملیات کپی را انجام دهد. برای این کار، ماشین $W$ می‌تواند به صورت رفت و برگشتی بین پیکربندی قدیم و پیکربندی در حال نوشتن جدید حرکت کند. با تطبیق موقعیت نسبی خانه‌ها نسبت به نمادهای $\#$، می‌تواند نماد صحیح بعدی را پیدا کرده و در انتهای نوار (که خالی است) بنویسد.
\end{itemize}

چون هر ماشین تورینگ استاندارد را می‌توان با یک \lr{WOTM} شبیه‌سازی کرد، هر ماشین استاندارد مانند $M$ ورودی $w$ را می‌پذیرد، اگر و تنها اگر ماشین شبیه‌ساز آن در قالب \lr{WOTM} همان رشته را بپذیرد. 

بنابراین، اگر ماشین تصمیم‌گیری برای $A_{WOTM}$ وجود داشته باشد، می‌توانیم از آن برای تصمیم‌گیری $A_{TM}$ استفاده کنیم که یک تناقض آشکار است. در نتیجه، این زبان تصمیم‌ناپذیر است.
